% !TEX root = ../algorithms.tex

\section{Дискретное преобразование Фурье (DFT), обратное DFT. Свёртки и умножение многочленов при помощи алгоритма FFT. Обобщения на другие кольца, NTT. Вычислительная схема FFT (преобразования бабочки).}

\subsection*{Свёртка и умножение многочленов}

\begin{Definition}{}{}
	Пусть заданы последовательности чисел
	$a_0,\dots,a_{n-1}$ и $b_0,\dots,b_{n-1}$.
	Определим последовательность $c_0,\dots,c_{2n-2}$ по формуле
	\[
	c_k=\sum_{t=0}^{n-1} a_t\, b_{k-t}, \qquad k=0,\dots,2n-2,
	\]
	(считаем $a_i=b_i=0$ при $i\notin[0,n-1]$).
	Эта операция называется \textbf{конволюцией (сверткой)} и обозначается $c=a*b$.
\end{Definition}

\subsubsection*{Умножение многочленов}

Рассмотрим многочлены
\[
A(x)=a_0+a_1x+\dots+a_{n-1}x^{n-1},
\qquad
B(x)=b_0+b_1x+\dots+b_{n-1}x^{n-1}.
\]
Тогда их произведение
\[
C(x)=A(x)B(x)= c_0+c_1x+\dots+c_{2n-2}x^{2n-2},
\]
где коэффициенты $c_k$ задаются конволюцией $c=a*b$. Таким образом, умножение многочленов и вычисление свёртки их коэффициентов — одна и та же задача.

\subsubsection*{Сложность}

\begin{itemize}
	\item Наивное умножение: $O(n^2)$.
	\item Карацуба: $O\!\left(n^{\log_2 3}\right)$.
	\item FFT (быстрое преобразование Фурье): $O(n\log n)$.
\end{itemize}

\subsection*{Идея FFT-умножения}

Хотим быстро вычислить свёртку (а значит, и произведение многочленов).
Для этого:

\begin{enumerate}
	\item Выбираем число точек $N$ так, чтобы $N\ge 2n-1$ и обычно $N=2^t$.
	Коэффициенты $A,B$ дополняем нулями до длины $N$.
	\item Быстро считаем значения в $N$ точках $A(\omega_k)$, $B(\omega_k)$, $k=0,\dots,N-1$,
	где $\omega$ — примитивный $N$-й корень из единицы.
	\item Перемножаем покомпонентно: $y_k = A(\omega_k)\,B(\omega_k)$.
	\item По значениям $y_k$ восстанавливаем коэффициенты $C$ обратным FFT.
\end{enumerate}

Ключевой факт: можно выбрать точки $\omega_k$ так, что и «вычисление значений»,
и «восстановление коэффициентов» делаются за $O(N\log N)$.

\subsection*{Реализация FFT}

Далее считаем, что $N=2^t$. Если исходная длина не степень двойки,
дополняем нулями.

\subsubsection*{Прямое FFT: из коэффициентов в значения}

Выбираем
\[
\omega = e^{2\pi i/N}=\cos\frac{2\pi}{N}+i\sin\frac{2\pi}{N},
\qquad
\omega_k=\omega^k.
\]

% -------------------------------------------------------
% Рисунок 1: N=8 корни из единицы (и квадраты → N=4)
% -------------------------------------------------------
\begin{figure}[H]
	\centering
	\begin{tikzpicture}[scale=1.5, >=Stealth, font=\small]
		
		% ---- Левый: N=8 ----
		\begin{scope}[xshift=-2.4cm]
			\draw[line width=0.9pt] (-1.3,0) -- (1.3,0) node[right,font=\small] {Re};
			\draw[line width=0.9pt] (0,-1.3) -- (0,1.33) node[above,font=\small] {Im};
			\draw[line width=0.8pt] (0,0) circle (1cm);
			
			\foreach \k in {0,...,7} {
				\pgfmathsetmacro\angr{\k*45}
				\pgfmathsetmacro\cx{cos(\angr)}
				\pgfmathsetmacro\cy{sin(\angr)}
				\draw[->, line width=1.0pt, blue!80!black] (0,0) -- (\cx,\cy);
			}
			
			\node[font=\scriptsize, right] at (1.08, 0.1) {$\omega_0$};
			\node[font=\scriptsize, above] at (0.2, 1.16) {$\omega_2$};
			\node[font=\scriptsize, left]  at (-1.08, 0.1) {$\omega_4$};
			\node[font=\scriptsize, below] at (0.2,-1.18) {$\omega_6$};
			
			\node[font=\scriptsize] at ({1.28*cos(45)},{1.28*sin(45)})   {$\omega_1$};
			\node[font=\scriptsize] at ({1.28*cos(135)},{1.28*sin(135)}) {$\omega_3$};
			\node[font=\scriptsize] at ({1.28*cos(225)},{1.28*sin(225)}) {$\omega_5$};
			\node[font=\scriptsize] at ({1.28*cos(315)},{1.28*sin(315)}) {$\omega_7$};
			
			\node[below,font=\footnotesize] at (0,-1.5) {$N=8$};
		\end{scope}
		
		% ---- Правый: квадраты → N/2=4 точки ----
		\begin{scope}[xshift=2.4cm]
			\draw[-, line width=0.9pt] (-1.3,0) -- (1.3,0) node[right,font=\small] {Re};
			\draw[-, line width=0.9pt] (0,-1.3) -- (0,1.3) node[above,font=\small] {Im};
			\draw[line width=0.8pt] (0,0) circle (1cm);
			
			\foreach \angr/\lab in {45/{$\omega_1$},135/{$\omega_3$},
				225/{$\omega_5$},315/{$\omega_7$}} {
				\draw[->, line width=0.9pt, dashed, gray!95!black] (0,0)
				-- ({cos(\angr)},{sin(\angr)});
				\node[font=\scriptsize, gray!95!black] at
				({1.28*cos(\angr)},{1.28*sin(\angr)}) {\lab};
			}
			
			\draw[->, line width=1.1pt, red!70](0,0)--(1,0);
			\node[font=\scriptsize,right] at (1.05, 0.2) {$(\omega_0)^2{=}\omega_0$};
			
			\draw[->, line width=1.1pt, red!70](0,0)--(0,1);
			\node[font=\scriptsize,above right] at (0.05,1.08) {$(\omega_1)^2{=}\omega_2$};
			
			\draw[->, line width=1.1pt, red!70](0,0)--(-1,0);
			\node[font=\scriptsize,left] at (-1.05,0.2) {$(\omega_2)^2{=}\omega_4$};
			
			\draw[->, line width=1.1pt, red!70](0,0)--(0,-1);
			\node[font=\scriptsize,below right] at (0.05,-1.12) {$(\omega_3)^2{=}\omega_6$};
			
			\node[below,font=\footnotesize,align=center] at (0,-1.5)
			{$\omega^{2k}$ — лишь $N/2{=}4$ точки};
		\end{scope}
		
	\end{tikzpicture}
	\caption{Слева — корни из единицы $\omega_k=e^{2\pi ik/8}$ при $N=8$.
		При возведении в квадрат восемь корней попарно сливаются: квадраты
		$\omega^{2k}$ (справа, красные стрелки) пробегают лишь $N/2=4$ различные точки.
		Именно это позволяет свести задачу размера $N$ к двум задачам размера $N/2$.
		Серые пунктирные стрелки — нечётные узлы $\omega_1,\omega_3,\omega_5,\omega_7$.}
	\label{fig:roots}
\end{figure}

Нужно вычислить
\[
y_k = A(\omega^k)=\sum_{j=0}^{N-1} a_j(\omega^k)^j,\qquad k=0,\dots,N-1.
\]

Разобьём $A$ на чётные и нечётные коэффициенты:
\[
A_0(x)=a_0+a_2x+\cdots+a_{N-2}x^{\frac{N}{2}-1},
\qquad
A_1(x)=a_1+a_3x+\cdots+a_{N-1}x^{\frac{N}{2}-1}.
\]
Тогда $A(x)=A_0(x^2)+xA_1(x^2)$, и при подстановке $x=\omega^k$ получаем
\[
A(\omega^k)=A_0(\omega^{2k})+\omega^k A_1(\omega^{2k}).
\]
Обозначим $y_k^{(0)}=A_0(\omega^{2k})$ и $y_k^{(1)}=A_1(\omega^{2k})$.
Так как значения $\omega^{2k}$ пробегают только $N/2$ различных точек,
достаточно посчитать $y_k^{(0)},y_k^{(1)}$ для $k=0,\dots,\frac{N}{2}-1$ (рекурсивно),
а затем «склеить» ответы:
\[
y_k = y_k^{(0)} + \omega^k\, y_k^{(1)},
\qquad
y_{k+\frac{N}{2}} = y_k^{(0)} - \omega^k\, y_k^{(1)}.
\]
(Здесь использовано $\omega^{k+N/2}=-\omega^k$.)

\subsubsection*{Оценка времени прямого FFT}

На каждом уровне рекурсии решаем $2$ подзадачи размера $N/2$
и делаем склейку за $O(N)$, поэтому $T(N)=2T\!\left(N/2\right)+O(N)=O(N\log N)$.

\subsubsection*{Обратное FFT: из значений в коэффициенты}

Прямое преобразование можно записать матрично:
\[
\begin{pmatrix}
	1 & 1 & 1 & \cdots & 1 \\
	1 & \omega^1 & (\omega^1)^2 & \cdots & (\omega^1)^{N-1} \\
	1 & \omega^2 & (\omega^2)^2 & \cdots & (\omega^2)^{N-1} \\
	\vdots & \vdots & \vdots & \ddots & \vdots \\
	1 & \omega^{N-1} & (\omega^{N-1})^2 & \cdots & (\omega^{N-1})^{N-1}
\end{pmatrix}
\begin{pmatrix}
	a_0\\ a_1\\ \vdots\\ a_{N-1}
\end{pmatrix}
=
\begin{pmatrix}
	y_0\\ y_1\\ \vdots\\ y_{N-1}
\end{pmatrix}.
\]
Это матрица Вандермонда $F_\omega$.

\begin{Statement}{}{}
	Выполнено $F_\omega \,F_{\omega^{-1}} = N\cdot \mathbf{I}$.
\end{Statement}

\begin{proof}
	Элемент произведения в позиции $(i,j)$ равен
	\[
	\sum_{k=0}^{N-1}(\omega^i)^k(\omega^{-j})^k
	=\sum_{k=0}^{N-1}(\omega^{i-j})^k.
	\]
	Если $i\neq j$, то $q:=\omega^{i-j}\neq 1$, и по формуле суммы геометрической прогрессии $\sum_{k=0}^{N-1} q^k=\frac{1-q^N}{1-q}$
	получаем
	\[
	\sum_{k=0}^{N-1}(\omega^{i-j})^k
	=\frac{1-(\omega^{i-j})^N}{1-\omega^{i-j}}
	=\frac{1-(\omega^N)^{\,i-j}}{1-\omega^{i-j}}
	=\frac{1-1}{1-\omega^{i-j}}=0.
	\]
	Если же $i=j$, то $\sum_{k=0}^{N-1}(\omega^{i-j})^k=\sum_{k=0}^{N-1}1=N$.
\end{proof}

Следовательно, обратное преобразование — это прямое FFT с корнем $\omega^{-1}$,
после чего нужно поделить все коэффициенты на $N$: $a = \frac{1}{N}\,F_{\omega^{-1}}\,y$.

\subsubsection*{Сложность умножения многочленов через FFT}

Нужно:
\begin{itemize}
	\item одно прямое FFT для $A$,
	\item одно прямое FFT для $B$,
	\item покомпонентное умножение ($O(N)$),
	\item одно обратное FFT.
\end{itemize}
Итого $O(N\log N)$, где $N$ — ближайшая степень двойки, не меньшая $2n-1$.
При этом $2n-1 \le N < 2(2n-1) < 4n$, то есть $N=\Theta(n)$, поэтому итоговая сложность равна $O(n\log n)$.

\subsection*{Обобщение: NTT (Number-Theoretic Transform)}

Вместо $\mathbb{C}$ можно работать в кольце/поле $R$, если:
\begin{itemize}
	\item в $R$ существует элемент $\omega$ порядка $N$ (примитивный $N$-й корень из единицы);
	\item число $N$ обратимо в $R$ (нужно для деления на $N$ в обратном преобразовании).
\end{itemize}

\subsubsection*{Пример (поле $\mathbb{Z}_p$)}

Пусть $R=\mathbb{Z}_p$, где $p$ — простое, и $N=2^k$, $p=c\cdot 2^k+1$.
Тогда $p-1$ кратно $N$, и существует первообразный корень $g$ по модулю $p$
(порождает мультипликативную группу порядка $p-1$).
Можно взять $\omega \equiv g^c \pmod p$; тогда $\omega$ имеет порядок $N$, и NTT работает полностью по модулю $p$.

\subsection*{Приложение: смысл коэффициентов DFT}

Пусть дан сигнал (последовательность) $a_0,a_1,\dots,a_{N-1}$.
Определим дискретное преобразование Фурье:
\[
y_k=\sum_{t=0}^{N-1} a_t\,\omega^{kt},\qquad k=0,\dots,N-1,
\quad \text{где }\ \omega=e^{2\pi i/N}.
\]
Так как $\omega^{kt}=\cos\!\left(\frac{2\pi k}{N}t\right)+i\sin\!\left(\frac{2\pi k}{N}t\right)$, то
\[
\Re(y_k)=\sum_{t=0}^{N-1} a_t\cos\!\left(\frac{2\pi k}{N}t\right),
\qquad
\Im(y_k)=\sum_{t=0}^{N-1} a_t\sin\!\left(\frac{2\pi k}{N}t\right).
\]

% -------------------------------------------------------
% Рисунок 2: Смысл DFT — выделение частоты
% -------------------------------------------------------
\begin{figure}[H]
	\centering
	\begin{tikzpicture}
		\begin{axis}[
			width=13cm, height=5.5cm,
			domain=0:1, samples=400,
			xmin=0, xmax=1,
			ymin=-1.25, ymax=1.25,
			axis lines=left,
			xlabel={$t/N$},
			xtick={0,0.25,0.5,0.75,1},
			xticklabels={$0$,$\tfrac{N}{4}$,$\tfrac{N}{2}$,$\tfrac{3N}{4}$,$N$},
			ytick={-1,0,1},
			tick label style={font=\small},
			label style={font=\small},
			legend style={at={(1.25,1.19)},anchor=north east,font=\small,
				draw=none, fill=white, fill opacity=0.7},
			]
			% Шумный сигнал: основная частота + гармоники
			\addplot[blue!70, thick, opacity=0.9]
			{ 0.55*sin(deg(2*pi*3*x))
				+ 0.20*sin(deg(2*pi*7*x + 0.5))
				+ 0.14*sin(deg(2*pi*11*x + 1.1))
				+ 0.10*sin(deg(2*pi*17*x + 0.3))
				+ 0.08*sin(deg(2*pi*23*x + 0.9))
			};
			\addlegendentry{сигнал $a_t$ (несколько гармоник)}
			
			% Доминирующая частота k=3
			\addplot[red!80, very thick, dashed]
			{ 0.55*sin(deg(2*pi*3*x)) };
			\addlegendentry{частота $k{=}3$: $|y_3|$ велик}
			
		\end{axis}
	\end{tikzpicture}
	\caption{Если в сигнале (синий) доминирует частота $k{=}3$,
		то коэффициент $y_3 = \sum_t a_t\,e^{2\pi i \cdot 3t/N}$ велик по модулю:
		вклады $a_t\cos\tfrac{2\pi\cdot 3\, t}{N}$ складываются (красный пунктир),
		а для других $k$ — взаимно компенсируются.}
	\label{fig:dft_meaning}
\end{figure}

Эти суммы — скалярные произведения сигнала с косинусом и синусом частоты $k$.
Если в сигнале сильно выражена гармоника частоты $k$, то вклады «складываются»,
и модуль $|y_k|$ получается большим; если такой частоты нет — вклады
в основном взаимно компенсируются и $|y_k|$ близок к нулю.

\subsection*{Вычислительная схема FFT: преобразование бабочки}

% ===========================================================
% Рис. 3. Одна «бабочка» FFT
% ===========================================================
\begin{figure}[H]\centering
	\begin{tikzpicture}[>=Stealth, font=\small, scale=0.95]
		\def\H{1.2}
		% входы
		\node[left] at (0,\H)  {$y_k^{(0)}$};
		\node[left] at (0,-\H) {$y_k^{(1)}$};
		% горизонтальные участки
		\draw (0,\H)  -- (1.5,\H);
		\draw (0,-\H) -- (0.5,-\H);
		% блок умножения
		\draw[rounded corners=3pt,thick] (0.5,-\H-0.45) rectangle (1.7,-\H+0.45);
		\node at (1.1,-\H) {$\times\omega^k$};
		% точки ветвления
		\filldraw (1.5,\H)  circle (2.5pt);
		\filldraw (1.7,-\H) circle (2.5pt);
		% «крест» бабочки
		\draw (1.5,\H)  -- (3.6,\H);
		\draw (1.7,-\H) -- (3.6,-\H);
		\draw (1.5,\H)  -- (3.6,-\H);
		\draw (1.7,-\H) -- (3.6,\H);
		% узлы суммирования
		\node[circle,draw,inner sep=1.5pt,fill=white,minimum size=7mm] (P) at (4.1,\H)  {$+$};
		\node[circle,draw,inner sep=1.5pt,fill=white,minimum size=7mm] (M) at (4.1,-\H) {$-$};
		\draw (3.6,\H)  -- (P);
		\draw (3.6,-\H) -- (M);
		% выходы
		\draw[->] (P) -- (6.1,\H)
		node[right] {$y_k \;=\; y_k^{(0)}+\omega^k\, y_k^{(1)}$};
		\draw[->] (M) -- (6.1,-\H)
		node[right] {$y_{k+n/2} = y_k^{(0)}-\omega^k\, y_k^{(1)}$};
	\end{tikzpicture}
	\caption{Одна «бабочка» — элементарный шаг схемы FFT.
		Нижний входной провод умножается на $\omega^k$;
		затем два провода попарно складываются и вычитаются.
		Знак минус объясняется тем, что $\omega^{k+n/2}=-\omega^k$.}
	\label{fig:fft-butterfly}
\end{figure}

Рекурсивная схема FFT состоит из двух блоков $\mathrm{FFT}_{n/2}$, к выходам которых добавляется один слой из $n/2$ бабочек: $k$-я бабочка умножает правый вход на $\omega^k$, а затем складывает и вычитает.

% ===========================================================
% Рис. 4. Рекурсивная схема FFT (два блока FFT_{n/2} + бабочки)
% ===========================================================
\begin{figure}[H]\centering
	\begin{tikzpicture}[>=Stealth, font=\small]
		\def\sep{0.62}
		% входы a_0..a_7 (0-based)
		\foreach \i in {0,...,7}{
			\pgfmathsetmacro{\x}{\i*\sep}
			\draw[thin] (\x,5.2) -- (\x,4.0);
			\node[above,font=\scriptsize] at (\x,5.2) {$a_{\i}$};
		}
		% чётные коэффициенты -> левый блок (позиции 0..3)
		\foreach \i/\t in {0/0,2/1,4/2,6/3}{
			\pgfmathsetmacro{\xs}{\i*\sep}
			\pgfmathsetmacro{\xt}{\t*\sep}
			\draw[thin] (\xs,4.0) -- (\xt,2.8);
		}
		% нечётные коэффициенты -> правый блок (позиции 4..7)
		\foreach \i/\t in {1/4,3/5,5/6,7/7}{
			\pgfmathsetmacro{\xs}{\i*\sep}
			\pgfmathsetmacro{\xt}{\t*\sep}
			\draw[thin,dashed] (\xs,4.0) -- (\xt,2.8);
		}
		
		\def\bw{2.0}
		\draw[thick,rounded corners=3pt] (0,1.8) rectangle (\bw,2.8);
		\node at (\bw/2, 2.3) {$\mathrm{FFT}_{n/2}$};
		\node[below,font=\scriptsize] at (\bw/2,1.8) {$y^{(0)}$ (чётные)};
		\pgfmathsetmacro{\rx}{4*\sep}
		\draw[thick,rounded corners=3pt] (\rx,1.8) rectangle (\rx+\bw,2.8);
		\node at (\rx+\bw/2, 2.3) {$\mathrm{FFT}_{n/2}$};
		\node[below,font=\scriptsize] at (\rx+\bw/2,1.8) {$y^{(1)}$ (нечётные)};
		
		\draw[thick,dashed,rounded corners=3pt]
		(-0.2,-0.1) rectangle (\rx+\bw+0.2,1.6);
		\pgfmathsetmacro{\midx}{(\rx+\bw)/2}
		\node[font=\small, align=center] at (\midx,0.75)
		{Слой из $n/2$ бабочек:\\
			$\times\omega^0,\;\times\omega^1,\;\dots,\;\times\omega^{n/2-1}$};
		
		\foreach \i in {0,...,3}{
			\pgfmathsetmacro{\x}{\i*\sep}
			\draw[thin] (\x,1.8) -- (\x,1.6);
			\pgfmathsetmacro{\xr}{\rx+\i*\sep}
			\draw[thin] (\xr,1.8) -- (\xr,1.6);
		}
		
		\foreach \i in {0,...,7}{
			\pgfmathsetmacro{\x}{\i*\sep}
			\draw[thin,->] (\x,-0.1) -- (\x,-0.8);
			\node[below,font=\scriptsize] at (\x,-0.8) {$y_{\i}$};
		}
	\end{tikzpicture}
	\caption{Рекурсивная схема FFT (на примере $n=8$): чётные коэффициенты
		$a_0,a_2,a_4,a_6$ поступают в левый блок $\mathrm{FFT}_{n/2}$ (он даёт $y^{(0)}$),
		нечётные $a_1,a_3,a_5,a_7$ — в правый (он даёт $y^{(1)}$);
		затем слой из $n/2$ бабочек объединяет результаты в $y_0,\dots,y_{n-1}$.}
	\label{fig:fft-recursive}
\end{figure}

\subsubsection*{Сложность}

Для схемы FFT глубина есть $O(\log n)$, а размер — $O(n\log n)$.
Глубина удовлетворяет рекуррентности $T(n)=T(n/2)+O(1)$, откуда $T(n)=O(\log n)$;
размер удовлетворяет $S(n)=2S(n/2)+O(n)$, откуда по мастер-теореме $S(n)=O(n\log n)$.

\begin{Proposition}{}{rem:fft-exercise}
	Пусть $y = F_{\omega} a$, где
	\[
	F_\omega = (\omega^{jk})_{j,k=0}^{n-1},
	\]
	а $\omega$ — примитивный корень степени $n$ из единицы.
	Тогда $F_\omega^2 = nP$,
	где $P$ — перестановка $P(a_0,a_1,\dots,a_{n-1}) = (a_0,a_{n-1},a_{n-2},\dots,a_1)$.
\end{Proposition}

\begin{proof}
	Для любого $a=(a_0,\dots,a_{n-1})^T$ и любого $i=0,\dots,n-1$ имеем
	\[
	(F_\omega^2 a)_i
	=
	\sum_{j=0}^{n-1}\omega^{ij}\sum_{k=0}^{n-1}\omega^{jk}a_k
	=
	\sum_{k=0}^{n-1} a_k \sum_{j=0}^{n-1}\omega^{j(i+k)}.
	\]
	Но
	\[
	\sum_{j=0}^{n-1}\omega^{jt}
	=
	\begin{cases}
		n, & t \equiv 0 \pmod n,\\
		0, & t \not\equiv 0 \pmod n.
	\end{cases}
	\]
	Следовательно, $(F_\omega^2 a)_i = n\,a_{-i \bmod n}$, то есть
	$F_\omega^2 a = n\,(a_0,a_{n-1},a_{n-2},\dots,a_1)^T = nPa$, откуда $F_\omega^2 = nP$.
	
	Отсюда, если $y=F_\omega a$ и $\tilde a = F_\omega y$, то
	\[
	\tilde a = F_\omega^2 a = nPa,
	\qquad
	a = \frac1n P\tilde a.
	\]
\end{proof}

Итак, чтобы восстановить $a$ по $y$, достаточно:
\begin{enumerate}
	\item вычислить $\tilde a = F_\omega y$;
	\item обратить порядок координат $\tilde a_1,\dots,\tilde a_{n-1}$, оставив $\tilde a_0$ на месте;
	\item разделить все координаты на $n$.
\end{enumerate}